\chapter{Grundlagen}
\label{chap:grundlagen}

\section{Netzwerkgrundlagen und DDoS-Angriffe}
\label{sec:netzwerk_ddos}
Moderne Kommunikationsnetze basieren auf der Übertragung digitaler Daten in standardisierten, paketbasierten Formaten zwischen Endsystemen. Dabei werden Informationen in Datenpakete zerlegt, mit Netzwerk- und Transportinformationen versehen und über Routing- und Switching-Verfahren an das Ziel weitergeleitet. Die Interoperabilität dieser Prozesse beruht auf Protokollen wie IP, TCP und UDP, deren Zusammenspiel die Zuverlässigkeit, Latenz und Bandbreitenverfügbarkeit des Datenaustauschs bestimmt. Ein Verständnis dieser Mechanismen ist essenziell, weil Angriffe auf Netzwerke typischerweise nicht auf die Vertraulichkeit von Daten, sondern auf die Verfügbarkeit und den Durchsatz der Kommunikationsinfrastruktur abzielen \autocite{mirkovicTaxonomyDDoSAttack2004}.

\subsection{Grundlagen der Datenkommunikation in Netzwerken}
Die Grundlage digitaler Kommunikation bildet das Schichtenmodell, insbesondere die Vermittlung auf der Netzwerk- und Transportschicht, bei der Adressen, Ports und Zustandsinformationen für den End-to-End-Austausch verwendet werden. In diesem Kontext ist die Bandbreite einer Verbindung ebenso relevant wie die Verarbeitungsleistung der beteiligten Server, Firewalls und Zwischenstationen, da jede dieser Komponenten potenziell durch Überlastung oder Fehlkonfiguration zu einem Engpass werden kann. Eine stabile Kommunikation setzt daher eine Balance zwischen Paketfluss, Verarbeitungszeit und Ressourcenallokation voraus, die im normalen Betrieb durch legitimen Nutzerverkehr geprägt ist und im Angriffsfall durch massiv gesteigerten, ungleichmäßigen oder aggressiv erzeugten Datenverkehr gestört wird.

\subsection{Definition, Ziele und Merkmale von DDoS-Angriffen}
Ein Distributed Denial of Service (DDoS)-Angriff bezeichnet die koordinierten, verteilten Versuche, die Verfügbarkeit einer Zielressource oder eines Netzwerksystems durch eine künstlich erzeugte Überlastung zu beeinträchtigen. Die Angreifer nutzen dabei in der Regel Botnets oder andere große Mengen kompromittierter Endgeräte, um eine hohe Last von mehreren Quellen gleichzeitig gegen ein Ziel zu senden und so die Wahrnehmung eines „normalen“ Verkehrs zu künstlich verzerren. Charakteristisch für solche Angriffe sind ihre Verteilung, ihre Skalierbarkeit, der Nutzungseffekt durch Amplifizierung und die Tendenz, legitime Anfragen durch Flutung mit pseudo-typischem Traffic zu verdrängen. Diese Merkmale erschweren nicht nur die Abwehr, sondern auch die präzise Unterscheidung zwischen normalem Spitzenverkehr und gezielter Störung, weshalb DDoS-Angriffe als besonders kritische Bedrohung für Internetdienste und kritische Infrastrukturen gelten.

\subsection{Typische DDoS-Angriffsvektoren und deren Auswirkungen}
Je nach Zielsystem und Protokoll können DDoS-Angriffe auf verschiedene Ressourcen wirken: So können volumetrische Angriffe die verfügbare Bandbreite des Netzes erschöpfen, protokollbasierte Angriffe die Verarbeitungslogik von Servern oder Load Balancern überlasten und anwendungsorientierte Angriffe gezielt einzelne Services wie Webanwendungen oder APIs schwächen. Die Folgen reichen von temporeller Service-Degradation und erhöhten Latenzzeiten bis hin zu vollständiger Erreichbarkeitslosigkeit, Betriebsunterbrechungen und wirtschaftlichen Schäden. Da die Effektivität eines DDoS-Angriffs häufig von der Kombination aus volumenstarker Flutung, Pseudo-Normalität des Verkehrs und einer hohen Verteilung der Angriffsquellen abhängt, muss die Abwehr in modernen Netzwerken sowohl auf Erkennung als auch auf adäquate Reaktionsmechanismen abzielen.

\section{Klassische Ansätze zur Angriffserkennung}
\label{sec:klassische_erkennung}
\subsection{Signaturbasierte und regelbasierte Erkennung}
Klassische Verfahren zur Angriffserkennung stützen sich primär auf signaturbasierte und regelbasierte Mechanismen. Bei der signaturbasierten Erkennung werden bekannte Muster bösartiger Aktivitäten, etwa charakteristische Bytefolgen, Header-Anomalien oder bekannte Ablaufsequenzen in Paketen, mit dem aktuellen Netzwerkverkehr abgeglichen. Regelbasierte Ansätze ergänzen dieses Prinzip durch explizit formulierte Heuristiken und Schwellwerte, beispielsweise zur Erkennung auffälliger Verbindungsraten, ungewöhnlicher Port-Nutzungen oder protokollwidriger Zustandsübergänge. Der zentrale Vorteil dieser Methoden liegt in ihrer hohen Nachvollziehbarkeit und ihrer geringen Fehlalarmrate bei bekannten Angriffstypen, da Entscheidungen auf klar definierten Kriterien beruhen und sicherheitsoperativ gut auditierbar sind.

\subsection{Anomaliebasierte Erkennung}
Neben diesen deterministischen Verfahren wurde die traditionelle Anomalieerkennung entwickelt, um auch bisher unbekannte Angriffe identifizieren zu können. Sie modelliert hierzu ein Referenzverhalten des legitimen Netzwerkverkehrs und markiert Abweichungen als potenziell sicherheitsrelevant. In klassischen Umsetzungen geschieht dies häufig über statistische Modelle, Zeitreihenanalysen oder einfache Distanzmaße im Merkmalsraum, etwa für Paketgrößen, Flussdauern oder Anfragefrequenzen. Solche Ansätze sind prinzipiell geeignet, neuartige oder leicht variierte Angriffsmuster sichtbar zu machen, stehen jedoch vor der Herausforderung, dass legitime Lastspitzen, saisonale Nutzungseffekte oder infrastrukturelle Änderungen ebenfalls als Abweichung erscheinen können.

\subsection{Grenzen klassischer Ansätze}
Die Grenzen klassischer Verfahren zeigen sich besonders deutlich bei verteilten und adaptiven Angriffen wie modernen DDoS-Kampagnen. Signatur- und regelbasierte Systeme reagieren oft verzögert auf neue Angriffsformen, da zunächst belastbare Muster extrahiert und Regeln manuell angepasst werden müssen. Gleichzeitig neigen traditionelle Anomalieansätze in dynamischen Produktionsnetzen zu erhöhten Fehlalarmen oder zu einer unzureichenden Sensitivität, wenn Angreifer ihren Traffic gezielt an normales Verhalten annähern. Hinzu kommen Skalierungsprobleme bei hohem Datenaufkommen sowie die begrenzte Fähigkeit, komplexe, mehrstufige Angriffsstrategien in Echtzeit robust zu erfassen. Diese Defizite bilden die wesentliche Motivation für datengetriebene und lernbasierte Erkennungsverfahren.

\section{KI-basierte Erkennung von Netzwerkangriffen}
\label{sec:ki_erkennung}
\subsection{Machine Learning für die Erkennung von Anomalien im Netzwerkverkehr}
Maschinelles Lernen ermöglicht die Erkennung von Anomalien im Netzwerkverkehr, indem Modelle aus historischen Daten Merkmalsmuster zwischen legitimem und bösartigem Verhalten lernen. Typische Eingabegrößen sind statistische Kennzahlen von Flüssen und Paketen, etwa Paketgrößen, Interarrival-Zeiten, Verbindungsdauern, Portnutzung und Volumenprofile, die als Merkmalsvektoren für Klassifikatoren wie Entscheidungsbäume, Random Forests, Support Vector Machines oder k-NN dienen. Im Vergleich zu rein regelbasierten Verfahren können diese Modelle nicht nur bekannte Angriffsmuster erkennen, sondern auch solche mit geringer Abweichung vom normalen Verkehr identifizieren, sofern sie in den Trainingsdaten ausreichend repräsentiert sind. Die Qualität des Ansatzes hängt dabei stark von Feature-Design, Datenqualität und der Robustheit gegenüber sich verändernden Netzwerkbedingungen ab.

\subsection{Deep Learning und neuronale Architekturen im Sicherheitskontext}
Deep Learning erweitert das maschinelle Lernen durch die automatische Extraktion hierarchischer Repräsentationen aus Roh- oder schwach vorverarbeiteten Netzwerkmerkmalen. Rekurrente Architekturen wie LSTM-Netzwerke eignen sich besonders für zeitliche Abhängigkeiten im Paketstrom, während Convolutional Neural Networks räumliche oder strukturelle Muster in aggregierten Verkehrsdaten erfassen können. Dadurch lassen sich komplexe, sequentielle Angriffsverläufe oft besser modellieren als mit klassischen, handgefertigten Regeln, insbesondere wenn Angriffe langsame, schrittweise Veränderungen im Netzwerkverhalten produzieren. Allerdings erfordern tiefere Modelle im Sicherheitskontext hohe Mengen an annotierten Daten, sorgfältige Regularisierung und eine aufwändige Hyperparameterwahl, um Generalisierung und Stabilität in Produktionsnetzen zu gewährleisten.

\subsection{Vorteile und Herausforderungen KI-basierter Ansätze für DDoS-Erkennung}
KI-basierte Ansätze bieten gegenüber traditionellen Verfahren den entscheidenden Vorteil, dass sie auf komplexe Muster im Verkehr reagieren und auch unbekannte oder adaptive Angriffsformen erkennen können, wenn diese im Trainingsbereich ausreichend repräsentiert sind. Sie eignen sich besonders für hochdimensionale, zeitlich variierende Ströme, wie sie bei volumetrischen DDoS-Angriffen und verteilter Aktivität auftreten, und können dadurch eine bessere Skalierbarkeit sowie eine höhere Erkennungsdynamik im Vergleich zu statischen Regeln erzielen. Gleichzeitig bleiben erhebliche Herausforderungen bestehen: Datenungleichgewichte, Konzeptdrift, Adversarial Attacks, hohe Rechenanforderungen und die Notwendigkeit einer robusten Bewertung auf realen Netzwerkdaten können zu Fehlalarmen oder verzögerter Reaktion führen. Für den praktischen Einsatz ist daher eine Kombination aus modellbasierter Erkennung, kontinuierlichem Monitoring und geeigneten Sicherheitsrichtlinien erforderlich, um die Aussagekraft der KI-Systeme mit der operativen Stabilität des Netzes in Einklang zu bringen.

\section{Machine Learning: Training, Evaluierung und Imbalanced Data}
\label{sec:ml_evaluierung}
Abschließend werden die mathematischen und methodischen Grundlagen des Modelltrainings, die Herausforderungen unbalancierter Datensätze und geeigneter Metriken sowie die Wirkung von Rauschen auf die Merkmale diskutiert.

\subsection{Grundlagen des maschinellen Lernens und der Modellbildung}
Die Modellbildung im maschinellen Lernen basiert auf der Approximation einer unbekannten Abbildung $f: \mathcal{X} \rightarrow \mathcal{Y}$ durch ein parametriertes Modell $\hat{f}_{\theta}$, dessen Parameter $\theta$ mittels empirischer Risikominimierung aus Trainingsdaten geschätzt werden. In überwachten Klassifikationsaufgaben erfolgt dies typischerweise durch die Minimierung einer Verlustfunktion wie der Kreuzentropie unter Verwendung gradientenbasierter Optimierung. Für eine belastbare Generalisierung auf unbekannte Daten sind eine saubere Trennung in Trainings-, Validierungs- und Testdaten, geeignete Regularisierungsverfahren sowie eine systematische Hyperparameteroptimierung erforderlich. Im Kontext der DDoS-Erkennung ist die Modellqualität dabei unmittelbar von der Repräsentativität zeitlich variierender Netzwerkzustände und der Stabilität der Merkmalsrepräsentation abhängig.

\subsection{Imbalanced Data und Klassenungleichgewicht}
Klassenungleichgewicht liegt vor, wenn die Verteilung der Zielklassen stark asymmetrisch ist, was in der Angriffserkennung häufig durch einen hohen Anteil legitimer gegenüber relativ seltenen bösartigen Verbindungen auftritt. Unter solchen Bedingungen können Lernverfahren bei Optimierung auf globale Genauigkeit eine systematische Mehrheitsklassen-Bias entwickeln und Minderheitsereignisse mit hoher sicherheitsrelevanter Bedeutung unzureichend erfassen. Zur Minderung dieses Effekts werden datenbasierte Strategien wie Over- und Undersampling, synthetische Datengenerierung (z. B. SMOTE), kosten-sensitive Verlustfunktionen oder klassengewichtete Lernziele eingesetzt. Eine methodisch korrekte Anwendung dieser Verfahren erfordert, dass Resampling ausschließlich auf den Trainingsdaten erfolgt, um Datenleckage und verzerrte Leistungsabschätzungen zu vermeiden.

\subsection{Evaluationsmetriken: Precision, Recall und F1-Score}
Für unausgewogene Klassifikationsprobleme sind Precision, Recall und F1-Score aussagekräftiger als die reine Accuracy, da sie Fehlklassifikationen klassenbezogen quantifizieren. Precision misst den Anteil korrekt erkannter Angriffe unter allen als Angriff klassifizierten Instanzen, Recall den Anteil erkannter Angriffe unter allen tatsächlich vorhandenen Angriffen, und der F1-Score kombiniert beide Größen als harmonisches Mittel $F1 = 2 \cdot \frac{\text{Precision} \cdot \text{Recall}}{\text{Precision} + \text{Recall}}$. Dadurch wird ein Modell nur dann hoch bewertet, wenn es sowohl eine geringe Fehlalarmrate als auch eine hohe Erkennungsrate erreicht. Für die DDoS-Detektion ist diese balancierte Betrachtung zentral, weil operative Sicherheitsanforderungen typischerweise einen Zielkonflikt zwischen Alarmpräzision und Angriffssensitivität adressieren.

\subsection{Merkmalsrauschen und Rauschsensitivität von Klassifikatoren}
\label{sec:merkmalsrauschen}
Rauschen in Trainings- und Testdaten tritt in zwei Formen auf. Klassenrauschen verfälscht das Label $c$, Attributrauschen die Merkmalswerte $\mathbf{x}$ selbst. Zhu und Wu haben beide Formen an 17 Datensätzen mit C4.5-Entscheidungsbäumen getrennt quantifiziert \autocite{zhuClassNoiseVs2004}. Für das Attributrauschen fanden sie einen Genauigkeitsverlust, der mit der Rauschrate wächst und umso stärker ausfällt, je enger das gestörte Merkmal mit der Zielklasse korreliert. In dieser Arbeit tritt nur Attributrauschen auf. Die Labels bleiben exakt, gestört wird ausschließlich die Beobachtung.

Formal entsteht die gestörte Beobachtung durch einen additiven Rauschterm
\begin{equation}
	\mathbf{y}=\mathbf{x}+\mathbf{n},\qquad \mathbf{n}\sim\mathcal{N}\!\left(\mathbf{0},\ \boldsymbol{\Sigma}_n\right),\quad \boldsymbol{\Sigma}_n=\operatorname{diag}\!\left(\sigma_{n,1}^2,\dots,\sigma_{n,d}^2\right),
\end{equation}
mit $\mathbf{n}$ unabhängig von $\mathbf{x}$ und $c$. Das ist das AWGN-Kanalmodell der Nachrichtentechnik mit Kanalfaktor $h=1$ \autocite{proakisDigitalCommunications2008}. Die Rauschleistung jeder Komponente legt das Signal-Rausch-Verhältnis fest, $\mathrm{SNR_{dB}}=10\log_{10}\!\left(P_j/\sigma_{n,j}^2\right)$ mit der Signalleistung $P_j$ des Merkmals $j$. Klassenbedingt gilt
\begin{equation}
	p(\mathbf{y}\mid c)=\bigl(p(\cdot\mid c)*\mathcal{N}(\mathbf{0},\boldsymbol{\Sigma}_n)\bigr)(\mathbf{y}).
\end{equation}
Die Faltung verbreitert jede Klassendichte. Die Überlappung nimmt zu, und mit ihr steigt der Bayes-Fehler, also die untere Schranke, die kein Klassifikator unterschreiten kann. Kein Modell bleibt unter AWGN verlustfrei. Robust bedeutet deshalb: Das Modell verliert nicht wesentlich mehr, als die gestiegene Klassenüberlappung erzwingt. Abschnitt~\ref{sec:klassenueberlappung} beziffert diesen unvermeidbaren Anteil für das verwendete Generatormodell.

Liegt das Rauschen auch im Training an, wirkt es zusätzlich auf die gelernte Abbildung. Bishop zeigt, dass Training mit additivem Eingangsrauschen kleiner Varianz $\eta^2$ bei quadratischer Fehlerfunktion einer Tikhonov-Regularisierung entspricht \autocite{bishopTrainingNoiseEquivalent1995}:
\begin{equation}
	\tilde{E}\approx E+\eta^{2}\,\Omega,\qquad \Omega=\frac{1}{2}\sum_{n}\sum_{k}\sum_{i}\left(\frac{\partial \hat{f}_k(\mathbf{x}_n)}{\partial x_{n,i}}\right)^{2}.
\end{equation}
$\Omega$ bestraft große Ableitungen der Modellausgänge $\hat{f}_k$ nach den Eingängen und glättet die Abbildung. Kleine Verschiebungen von $\mathbf{x}$ ändern die Ausgabe eines so trainierten Modells weniger stark. Die Herleitung setzt kleines $\eta^2$ und eine differenzierbare Abbildung voraus. Für das Q-Netz, dessen Huber-Verlust um den Nullpunkt quadratisch verläuft (Abschnitt~\ref{sec:dqn_architektur}), und für das MLP gilt der Zusammenhang näherungsweise, für den Random Forest als stückweise konstantes Modell nicht. Bei 5\,dB erreicht die Rauschstandardabweichung im hier verwendeten Modell $0{,}31$ und übersteigt die Klassenstreuung $\sigma=0{,}20$ (Abschnitt~\ref{sec:awgn_kanal}). Dort taugt der Regularisierungsblick nur noch als qualitative Erklärung. Er begründet trotzdem, warum alle Modelle dieser Arbeit mit pro Fluss gemischtem SNR trainiert werden: Ein rein auf sauberen Merkmalen trainiertes Modell träfe im Test auf eine Verteilung, die es nie gesehen hat.

Klassifikatoren mit unterschiedlicher Ausgangsleistung lassen sich über mehrere Rauschstufen nur fair vergleichen, wenn der Leistungsverlust auf die Ausgangsleistung bezogen wird. Ein Modell, das schon ohne Rauschen schwach ist, kann sonst allein wegen seines geringen Spielraums nach unten robust erscheinen. Sáez et al. führen dafür die Equalized Loss of Accuracy (ELA) ein \autocite{saezEvaluatingClassifierBehavior2016}. Abschnitt~\ref{sec:eval_metriken} überträgt dieses Maß auf den F1-Score.